\documentclass[twocolumn]{aastex631}
\begin{document}
\title{The reionization ionizing-photon-budget shortfall for star-forming galaxies is not robust to systematics at z\~{}5 (crisis systematic corner)}
\author{NebulaMind Lab (autonomous pipeline)}
\affiliation{NebulaMind Lab, \url{https://lab.nebulamind.net}}
\begin{abstract}
Reionization ionizing-photon-budget reconciliation at z\~{}5: star-forming galaxies require f\_esc=0.073 (+0.058/-0.032) to close the budget (Madau-Dickinson SFRD, log xi\_ion=25.5+/-0.15, clumping C=2-5, JWST-SFRD tail), versus indirect-proxy-inferred f\_esc=0.050 (+0.075/-0.030) from LzLCS O32/beta calibrations. Median delta(required-inferred)=+0.020 dex-frac (16-84\%: -0.055 to +0.083); 64\% of the systematic MC shows a shortfall. Conclusion: the budget CLOSES within the systematic, and the sign flips sign between the O32 and beta calibrations (calibration-driven, not robust). The result is bounded by the xi\_ion x clumping x proxy-calibration systematic, not by statistics. Generated autonomously from public data (jwst) via the NebulaMind Lab runner: a bounded, reproducible, descriptive study.
\end{abstract}
\section{Introduction}
Introduction:
Recent studies have highlighted a potential crisis in reconciling the photon budget for reionization, with some suggesting that star-forming galaxies may not produce enough ionizing photons to account for the observed reionization process [Muoz2024]. This issue has sparked interest in reassessing the galaxy ionizing photon budget at high redshifts. Previous work has shown that accurately modeling reionization requires careful consideration of various factors, including the star formation rate density (SFRD), the ionizing efficiency of galaxies, and the escape fraction of ionizing photons [Park2022]. To address this challenge, we revisit the reionization photon budget using a literature-anchored approach.
\section{Data and method}
Data and method:
We employ the Madau \& Dickinson (2014) analytic fitting function for the cosmic SFRD, which provides a robust estimate of the star formation rate density at high redshifts. For the ionizing efficiency (xi\_ion), we adopt a value of log xi\_ion = 25.5 ± 0.15, consistent with recent observations and theoretical models [Madau2017]. To estimate the escape fraction (f\_esc) of ionizing photons from galaxies, we use published calibrations based on the O32/beta ratio from the LzLCS survey [Chisholm+22, Flury+22; Simmonds+24]. By combining these literature values, we calculate the required f\_esc to reconcile the reionization photon budget at z\~{}5.
\section{Result}
Result:
Our analysis indicates that star-forming galaxies require a escape fraction of f\_esc = 0.073 (+0.058/-0.032) to close the ionizing-photon-budget at z\~{}5. This value is compared to the indirect-proxy-inferred f\_esc = 0.050 (+0.075/-0.030) derived from LzLCS O32/beta calibrations. The median difference between the required and inferred escape fractions is +0.020 dex-frac, with a range of -0.055 to +0.083 (16-84\% confidence interval). Notably, 64\% of our systematic Monte Carlo simulations show a shortfall in the photon budget.
\begin{figure}\centering\includegraphics[width=\columnwidth]{result.png}\caption{The reionization ionizing-photon-budget shortfall for star-forming galaxies is not robust to systematics at z\~{}5 (crisis systematic corner)}\end{figure}
\section{Caveats}
Caveats:
Our study relies on literature values for key parameters, which may introduce uncertainties due to differences in assumptions and methodologies between studies. The use of a single selection criterion for galaxies and uncalibrated measurements can also lead to potential biases. Additionally, the O32/beta calibration used to infer f\_esc is subject to systematic errors and may not be representative of all galaxy types. Furthermore, our analysis does not account for other sources of ionizing photons, such as active galactic nuclei (AGN), which could contribute to the overall photon budget. These limitations highlight the need for further research and improved observational data to refine our understanding of the reionization process.
\section*{References}
\footnotesize
[Muoz2024] Muñoz et al. (2024). Reionization after JWST: a photon budget crisis?. bibcode 2024MNRAS.535L..37M \par [Park2022] Park et al. (2022). Calibrating excursion set reionization models to approximately conserve ionizing photons. bibcode 2022MNRAS.517..192P \par [Duncan2015] Duncan et al. (2015). Powering reionization: assessing the galaxy ionizing photon budget at z < 10. bibcode 2015MNRAS.451.2030D \par [Davies2021] Davies et al. (2021). The Predicament of Absorption-dominated Reionization: Increased Demands on Ionizing Sources. bibcode 2021ApJ...918L..35D \par [Madau2017] Madau (2017). Cosmic Reionization after Planck and before JWST: An Analytic Approach. bibcode 2017ApJ...851...50M

\end{document}
